% ===========================================================================
\chapter{Setting up}\label{ch:setup}

Make a folder for the project and go into it. Every file in this tutorial
goes in this one folder.

\begin{shell}
mkdir agent
cd agent
\end{shell}

Copy two folders into it from \code{examples/coding_agent} of the Nex
distribution:

\begin{itemize}
\item \code{sample_project}, the project the agent will work on: a small
  lending library, with a bug;
\item \code{sessions}, which holds a recorded session, for
  Chapter~\ref{ch:replay}.
\end{itemize}

\section{The sample project}

\code{sample_project} has three Nex files and a \code{README.md}. A
\code{Loan} is a book lent to someone on some day (days are numbered from
the day the library opened), and is due back 14 days later:

\begin{excerpt}{sample_project/loans.nex}
  -- The last day the book may be returned without a fine.
  due_on(): Integer do
    result := lent_on + LOAN_DAYS
  end

  -- How many days late a book returned on `day` is: 0 if on time.
  days_late(day: Integer): Integer
  require
    not_before_loan: day >= lent_on
  do
    if day >= due_on() then
      result := day - due_on() + 1
    end
  ensure
    never_negative: result >= 0
  end
\end{excerpt}

A \code{Fine_Policy} charges so many cents for each day late, up to a
maximum. And \code{checks.nex} checks both. Run the checks:

\begin{shell}
cd sample_project
nex checks.nex
\end{shell}

Four of the seven fail. You may already see why, but don't fix it: that is
the agent's job. (If you do fix it, copy the folder again.)

\section{A first file: text helpers}

An agent spends most of its time moving text around: file contents,
command output, the model's replies. The first file holds the small text
utilities the rest of the project uses.

First, a word about printing. Nex's \code{print} shows a string with its
quotes, as \code{"hello"}, which is handy when debugging but not in a
program's output. Every Nex program has a \code{Console} called
\code{console}, ready to use without creating it, and
\code{console.print_line("hello")} prints the text as it is, on a line of
its own. We use it for all our output.

Nex strings have \code{index_of}, which finds the first place a piece of
text occurs. We will often need the next place after some position, to
count occurrences or to split text into lines, so the first helper does
that. It searches the rest of the string, from \code{start} on, and adds
\code{start} back to the position it finds:

\begin{newfile}{helpers.nex}
-- helpers.nex: small text utilities the rest of the agent uses, and
-- finding the user's files.

-- Where `part` first occurs in `text` at or after position `start`, or -1.
function index_from(text, part: String, start: Integer): Integer
require
  part_given: part /= ""
  start_in_range: start >= 0 and start <= text.length
do
  let at: Integer := text.substring(start, text.length).index_of(part)
  result := when at < 0 then -1 else start + at end
ensure
  not_found_or_after: result = -1 or result >= start
end
\end{newfile}

The \code{require} clause is the function's precondition: what must be true
when it is called. Searching for the empty string makes no sense, and a
start outside the string is a mistake by the caller. If a call breaks a
precondition, Nex stops with an error that names the clause, such as
\code{Precondition violation: part_given}, instead of carrying on with
nonsense. The \code{ensure} clause, the postcondition, is what the function
promises in return. We will lean on contracts like these throughout: they
say what each piece of code expects and guarantees, and Nex checks them as
the program runs.

\code{when} \ldots\ \code{then} \ldots\ \code{else} \ldots\ \code{end} is
an \code{if} that gives a value.

Now we can count how many times a piece of text occurs. This will be the
heart of the agent's safe editing in Chapter~\ref{ch:editing}: an edit is
only allowed if the text to replace occurs exactly once.

\begin{append}{helpers.nex}
-- How many times `part` occurs in `text`, not counting overlaps:
-- count_occurrences("banana", "an") is 2.
function count_occurrences(text, part: String): Integer
require
  part_given: part /= ""
do
  let at: Integer := index_from(text, part, 0)
  from until at < 0 do
    result := result + 1
    at := index_from(text, part, at + part.length)
  end
ensure
  never_negative: result >= 0
end

-- `text` with its first `before` replaced by `after`.
function replace_first(text, before, after: String): String
require
  occurs: text.index_of(before) >= 0
do
  let at: Integer := text.index_of(before)
  result := text.substring(0, at) + after + text.substring(at + before.length, text.length)
end
\end{append}

\code{from} \ldots\ \code{until} \ldots\ \code{do} \ldots\ \code{end} is
Nex's loop: it repeats the body until the condition becomes true. An
\code{Integer} \code{result} starts at 0, so \code{count_occurrences} can
add to it straight away.

Why not use the strings' own \code{split} to cut text into lines? Because
it takes a regular expression rather than plain text, and it drops empty
pieces at the end. Files are made of lines, and an agent has to show them
faithfully, blank lines and all, so we write our own:

\begin{append}{helpers.nex}
-- The lines of `text`, without their line breaks. A final line break
-- does not start another line: lines_of("a\nb\n") is ["a", "b"], and
-- lines_of("") is [].
function lines_of(text: String): Array[String] do
  let start: Integer := 0
  from until start >= text.length do
    let stop: Integer := index_from(text, "\n", start)
    if stop < 0 then
      stop := text.length
    end
    result.add(text.substring(start, stop))
    start := stop + 1
  end
end

-- Lines joined back together, each ending in a line break.
function joined_lines(lines: Array[String]): String do
  result := ""
  across lines as line do
    result := result + line + "\n"
  end
end
\end{append}

An \code{Array} \code{result} starts out empty, ready to \code{add} to.
\code{across} \ldots\ \code{as} \ldots\ \code{do} runs its body once for
every element.

When the agent shows the model a file, it numbers the lines, so that the
model can say ``line 32'' and we can see what it means. \code{pad_start}
adds spaces on the left until the number is 5 characters wide:

\begin{append}{helpers.nex}
-- `lines` numbered from `first`, the way read_file shows a file:
--    12  some code
function numbered(lines: Array[String], first: Integer): String do
  result := ""
  let n: Integer := first
  across lines as line do
    result := result + n.to_string.pad_start(" ", 5) + "  " + line + "\n"
    n := n + 1
  end
end
\end{append}

\section{Keeping output short}

A model reads everything we send it, and pays for it, in time and money:
models read text as \emph{tokens}, pieces of a few characters each, and
are priced by the token. A command can print a megabyte. So every tool's
output will be cut down to a sensible size before the model sees it.
Which part to keep? Usually the beginning and the end: the beginning says
what a command was doing, the end how it finished (the error, or the
summary line of a test run). So \code{shortened} drops the middle, and says
how much it dropped:

\begin{append}{helpers.nex}
-- `text` cut down to at most about `limit` characters by dropping the
-- middle. The start and the end of a long output are usually the useful
-- parts: what a command was doing, and how it finished.
function shortened(text: String, limit: Integer): String
require
  sensible_limit: limit >= 100
do
  if text.length <= limit then
    result := text
  else
    let keep: Integer := (limit - 60) / 2
    let cut: Integer := text.length - 2 * keep
    result := text.substring(0, keep)
      + "\n[... " + cut.to_string + " characters cut ...]\n"
      + text.substring(text.length - keep, text.length)
  end
ensure
  short_enough: result.length <= limit
end
\end{append}

The 60 characters left over are room for the note in the middle; the
postcondition checks that it fits.

The next helpers are for what the agent prints for you: a one-line preview
of a long text, counts like ``3 lines'' and ``2 replies'', and indenting:

\begin{append}{helpers.nex}
-- The first line of `text`, cut to `limit` characters: a one-line preview.
function preview(text: String, limit: Integer): String
require
  sensible_limit: limit >= 10
do
  result := text.trim()
  let stop: Integer := result.index_of("\n")
  if stop >= 0 then
    result := result.substring(0, stop) + " ..."
  end
  if result.length > limit then
    result := result.substring(0, limit - 3) + "..."
  end
ensure
  short_enough: result.length <= limit
end

-- "1 line", "3 lines", "2 replies", "4 matches".
function count_text(n: Integer, thing: String): String do
  result := n.to_string + " " + (when n = 1 then thing else plural(thing) end)
end

-- The plural of a noun, for the usual cases: lines, days, replies, matches.
function plural(noun: String): String do
  let ending: String := noun.substring((noun.length - 2).max(0), noun.length)
  if noun.ends_with("y") and not "aeiou".contains(ending.substring(0, 1)) then
    result := noun.substring(0, noun.length - 1) + "ies"
  elseif noun.ends_with("s") or noun.ends_with("x") or ending = "ch" or ending = "sh" then
    result := noun + "es"
  else
    result := noun + "s"
  end
end

-- Every line of `text` with `prefix` in front.
function indented(text, prefix: String): String do
  result := ""
  across lines_of(text) as line do
    result := result + prefix + line + "\n"
  end
end
\end{append}

The last helper finds a file the user names. The \code{nex} launcher runs
programs from its own installation folder, so a plain name like
\code{sample_project} would be looked for in the wrong place. Every program
has \code{program}, the running \code{Program}, which knows the folder you
ran \code{nex} from: \code{program.working_directory}. We put that in front
of any name that isn't already a full path:

\begin{append}{helpers.nex}
-- The full path of a file the user named. The `nex` launcher runs
-- programs from its own installation folder, so a relative name is
-- resolved against the folder the program was started from.
function user_path(raw: String): String do
  if raw.starts_with("/") then
    result := raw
  else
    result := program.working_directory + "/" + raw
  end
end
\end{append}

\section{Running a try-it script}

Let's check that it works. Make this file and run it:

\begin{tryit}{try_helpers.nex}
intern Helpers

console.print_line(count_occurrences("banana", "an").to_string)
console.print_line(lines_of("one\n\nthree\n").to_string)
console.print_line(numbered(["class Loan", "create"], 9))
console.print_line(preview("exit code 1\nPASS a loan is due ...", 30))
console.print_line(count_text(1, "line") + ", " + count_text(2, "reply") + ", "
  + count_text(3, "match"))
let long: String := ""
repeat 30 do
  long := long + "0123456789"
end
console.print_line(shortened(long, 160))
\end{tryit}
\tryoutput{try_helpers}

\code{intern Helpers} finds \code{helpers.nex} in the same folder and loads
everything in it. The file does not need a class called \code{Helpers}:
ours holds only plain functions. \code{repeat 30 do} runs its body 30
times.

Try-it scripts are not part of the project. Keep them around to experiment
with, or delete them when you are done.
